\section{Background}
Mechanistic interpretability currently can do supervised steering well
or unsupervised steering poorly. The Google DeepMind mechanistic
interpretability team has pivoted from ambitious reverse-engineering to
pragmatic interpretability judged on proxy tasks, and has deprioritized
SAE research: SAE probes did not generalize better out of distribution,
and SAEs did not beat a steering vector built from one contrastive pair
of prompts at suppressing a model's awareness of being
evaluated~\cite{gdmprogress2025}. On current benchmarks the supervised
difference of means steers better than SAE latents, though worse than
prompting~\cite{wu2025axbenchsteeringllmssimple}. But a supervised
method finds only the concepts an experimenter already has labels for,
and a direction that changes outputs as intended need not be one the
model computes with: a subspace intervention optimized for its effect
can instead activate a dormant pathway~\cite{makelov2023}. The dominant
unsupervised paradigm is the sparse autoencoder (SAE). A model's hidden
state is approximated as a finite state machine consisting of discrete
features. The first implementations used a $\relu$ activation and
sparsity loss~\cite{cunningham2023,bricken2023}, though more recent
implementations have used a $\topk$ activation function without a
sparsity loss~\cite{gao2024}.

\subsection{Common problems with SAEs}
The motivation is to decompile a model to a finite state machine (FSM).
The limitation is that such an FSM is only as interpretable as an
enumeration of all possible states. Even a mixture of two structureless
processes, one constant and one a fair coin, needs infinitely many
states to track which component is generating, so the structure lives
in how the components overlap~\cite{crutchfield2026}.

\subsubsection{Features are not disentangled}
A representation is disentangled when a change in one factor of variation
changes one factor of the representation~\cite{locatello2019}. A
persistent problem with SAEs is feature absorption~\cite{chanin2025}:
when one feature implies another, sparsity is better served by folding
the general feature's direction into the specific feature's latent and
switching the general latent off there, so a latent that looks
monosemantic silently fails to fire on some inputs that have its feature
(a ``starts with S'' latent that is silent on ``short''). A related
failure, feature composition, learns co-occurring concepts as one latent.
Matryoshka SAEs~\cite{bussmann2025} reduce both by training nested
prefixes of a single dictionary to reconstruct the input on their own,
so early latents learn general features and later ones specific
features. Tree SAEs~\cite{cao2026} make the hierarchy explicit by gating
each latent on an assigned parent.

\subsubsection{Features are not cleanly causal}
Sparsity trades off against reconstruction fidelity~\cite{gao2024}.
Steering with an SAE feature is rarely clean: adding its decoder
direction moves the logits differently in different contexts and
disturbs unrelated downstream features. How much varies by feature and is
partly predictable before intervening, from the crowding of its decoder
direction and its direct effect on the logits, though which predictor
works differs between models and dictionaries~\cite{duan2026}.

\subsubsection{Features are not closed-form}
A major concern is how well an interpretability technique generalizes
out-of-distribution. Ideally we would like to infer behavior by looking
solely at model weights. Bilinear MLPs make this partly possible: the
interaction matrix of any output direction can be eigendecomposed
without data, although in deep models the output directions to
decompose still come from data-dependent SAEs~\cite{pearce2025}.

\subsubsection{Forward passes are expensive}
Efficiency is also a problem. Although the sparsity constraint means
that only a handful of features are active for a given input, the
entire feature space needs to be computed for each
sample~\cite{gao2024,mudide2025}. Mixture-of-experts (MoE) SAE variants
reduce compute costs by dividing feature space between smaller
submodules. Examples include Switch SAEs, which route each input to one
expert~\cite{mudide2025}, and Scale SAEs, which activate several to curb
features duplicated across
experts~\cite{xu2025redundancydiversespecializedmultiexpert}.

Less explored is the effect of the structure imposed by the MoEfied
architecture on the interpretability of feature space. A common approach
to scalably interpreting features is to have an LLM explain each feature
from its activating examples and score the explanation by how well an
LLM judge, given only the explanation, separates activating from
non-activating contexts~\cite{paulo2025}; in the results here one model
both writes and judges the explanations. With an unstructured feature
space, it quickly becomes intractable to interpret all permutations of
features this way. It is however possible to instead apply this
technique to expert routings. The ReLU feed-forward layers of a trained
transformer can themselves be split into experts of co-activating
neurons, keeping over 95\% of performance with 10--30\% of the
parameters~\cite{zhang2022}. Language models appear to represent a
categorical concept as a simplex whose vertices are the vectors of its
values, with hierarchically related concepts in orthogonal subspaces, so
that a hierarchy is a direct sum of simplices; differences between value
vectors are the steering vectors of binary contrasts~\cite{park2025}. We
conjecture that MoEfication into classifiers can reveal such latent
contrastive steering vectors, since an entry swap writes exactly such a
difference. Penalizing the cosine between each SAE decoder direction and
its nearest neighbor cuts measured feature absorption by about two
thirds relative to BatchTopK SAEs, at slightly lower explained variance
and unchanged autointerp scores~\cite{korznikov2025}; that penalty is
unstructured, where the Ontologizer asks for orthogonality between heads
and contrast within them.

Existing attempts to route activations to expert SAEs focus on routing
similar inputs to the same expert~\cite{mudide2025}, and score an expert
as specialized when its features resemble each other and not other
experts'~\cite{xu2025redundancydiversespecializedmultiexpert}.
Contrastive learning wants the
opposite: training on things at opposite ends of a shared axis. We don't
just want to find similar features, we want to find contrastive vectors
in feature space that correspond to what a model considers a natural
category.

\subsection{Disjoint support}
\label{sec:bg-disjoint}
Concatenated heads give each head its own block of the
dictionary's coordinates, so the atoms of different heads have disjoint
supports by construction. Disjoint support is the strong form of
orthogonality: for non-negative atoms the two coincide
(Equation~\ref{eq:disjoint}), while signed vectors can be orthogonal by
cancellation and share every coordinate. Summed heads can only be asked
for it, and they do not deliver it: no trained summed model partitions its
coordinates between heads, and penalizing the overlap $\omega$ on GPT-2
produced a partition only by moving atom energy into directions the
decoder discards or shrinks, where reconstruction cannot see it
(Sections~\ref{sec:disjoint} and~\ref{sec:prelim-orth}). Disjoint support
bears on each of the four problems above, though unequally. In each case
the guarantee holds between the heads of one layer and in the dictionary;
what reaches the output passes through a shared decoder, which in general
maps the heads' blocks onto overlapping subspaces.

\paragraph{Disentanglement.}
Absorption needs two things: a general latent that switches off where a
specific one fires, and a specific latent whose direction takes over the
general one's. The head structure removes the first, since every head
writes on every input, choosing among its entries, and so cannot fall
silent where its feature applies. Disjoint support removes the second: no
entry of one head can carry a component of another head's block. Feature
composition is discouraged rather than excluded. Heads already make the
code factorial, $k^h$ codes from $hk$ atoms, so concepts that vary
independently can take separate heads, and under a quadratic cost on
activity and a bound on the weights~\cite{whittington2023,dorrell2025}, a
non-negative code gives two variables disjoint supports when their joint
range is spread enough, roughly when every combination of their extreme
values occurs~\cite{dorrell2025}. On SONAR the data meet that condition:
of the entry pairs of two heads that independence predicts at least ten
times, at most 0.24\% never occur in any layer of a SONAR model. The
summed models carry no such cost and do not partition; concatenation
imposes the partition by construction. Nothing yet forces a head to hold
one concept rather than a conjunction of two, nor can any method that sees
only the marginal data recover its factors without matching inductive
biases~\cite{locatello2019}. The heads are such a bias, and
Section~\ref{sec:prelim-seed} measures how far it reaches.

\paragraph{Causality.}
An SAE's steering is unclean partly because its decoder directions crowd
one another~\cite{duan2026}. Between heads, nothing crowds in the
dictionary. Forcing a head's entry rewrites only that head's block of the
layer's output, orthogonal to everything the other heads of the layer
write, so interventions on different heads move the output along
orthogonal directions and each head's contribution can be read off its
own block. Summed heads also compose additively, by linearity, but their
writes overlap, so one head's contribution cannot be read apart from the
others'. What disjoint support does not govern is how the other heads
read a stepped input: measured by re-encoding, an input-space step moves
as many other heads as a random step of its length, because the heads sit
almost on their decision boundaries and downstream layers reclassify a
changed residual (Section~\ref{sec:prelim-depth}). A disjoint write
removes neither cause.

\paragraph{Closed form.}
A static reading of the dictionary, which interprets a head from its
atoms without data, needs heads that own disjoint coordinates as a
property of the weights, for every input rather than on the data measured
(Section~\ref{sec:disjoint}). In the summed models that premise fails:
each head's usage covers 52--98\% of the coordinates. Concatenated heads
supply it, and when each head's decoder block has full column rank, every
part of a head's atoms reaches the output, unlike the partitions the
overlap penalty produced. It is a premise, not yet a reading: decoding
atoms through the paired decoder has not described what an entry does in
the models scored so far (Section~\ref{sec:prelim-autointerp}).

\paragraph{Cost.}
Disjoint support does least here. Concatenation divides the dictionary
among the heads, so at matched width $e$ each head stores $k$ atoms of
$e/h$ coordinates rather than $e$, $h$ times fewer dictionary parameters,
and a hard code's dictionary step is $h$ lookups that together write one
$e$-vector. But the classifier still scores all $hk$ entries for every
input, as an SAE scores all of its latents, so the forward pass is no
cheaper on the encoder side. Narrow blocks also cost reconstruction: on
GPT-2 layer-8 residuals at $d_{\mathrm{head}}{=}k{=}32$ under non-negative
atoms and constant temperature, where orthogonal atoms within a head need
one coordinate each, concatenated heads had the worst held-out error of
the orthant arms ($\mathrm{FVU}_w$ $0.416$ against $0.154$--$0.293$;
working notes), whereas giving each head headroom
($d_{\mathrm{head}}{=}48 > k{=}32$) and annealing the temperature brings
concatenated heads to $\mathrm{FVU}$ $0.186$ ($0.164$ at $k{=}64$), ahead
of direct atoms' $0.195$ (Section~\ref{sec:prelim-gpt2}).


